% GENERATED FILE. Edit canonical lessons, then run npm run generate:books.
% canonical-source-hash: fnv1a64:cbd4804855d09886
% canonical-lessons: GU-C312-patrakar, GU-C312-ghatna, GU-C312-kisso, GU-C312-aheval, GU-C312-varnavvu

\chapter{The Reporter}
\label{ch:gu-the-reporter}
\hlchaptermodality{312}

\begin{chapteropening}
\textbf{By the end of this chapter:} \emph{I can say what a reporter asked and wrote.}

Complete the last lesson of chapter 312: I can say what a reporter asked and wrote.
\end{chapteropening}

\section[\texorpdfstring{\gu{પત્રકાર} (patrakār)}{પત્રકાર (patrakār)}]{\gu{પત્રકાર} (patrakār) — a journalist}

\label{lesson:GU-C312-patrakar}



\begin{quote}
Before the new one: say the Gujarati for a radio, then the Gujarati for to inform.
\end{quote}

\subsection*{You'll want to know: \gu{પત્રકાર}}
\textbf{\gu{પત્રકાર}} — \emph{patrakār} — \textquotedblleft{}a journalist\textquotedblright{}.

\begin{quote}
\gu{એક} \gu{પત્રકારે} \gu{મીરાને} \gu{ઘણા} \gu{પ્રશ્ન} \gu{પૂછ્યા}.
\end{quote}

\begin{quote}
\emph{ek patrakāre Mīrāne ghaṇā praśna pūchyā.}
\end{quote}

\begin{quote}
A journalist asked Mira a lot of questions.
\end{quote}

A masculine or feminine noun, by the person meant.

\begin{cousinweb}
\textbf{\gu{પત્રકાર}} (\emph{patrakār}) is built from \textbf{\gu{પત્ર}}, a letter or a printed sheet, from Sanskrit \emph{patra}, \textquotedblleft{}a leaf\textquotedblright{}, and \emph{-kār}, \textquotedblleft{}a maker\textquotedblright{} (Monier-Williams): the one who makes the paper. It is used of a man or a woman.
\end{cousinweb}

\begin{grammarlens}[title={What you've built}]
One more word for telling what happened.
\end{grammarlens}

\subsection*{Your turn}
\begin{itemize}
\raggedright
  \item \emph{Say it:} \emph{patrakār}
  \item \emph{Say it:} \emph{patrakār}, once more
  \item \emph{Say it:} \emph{ek patrakāre Mīrāne ghaṇā praśna pūchyā.} — A journalist asked Mira a lot of questions.
  \item \emph{Recall:} say the Gujarati for a radio, then the Gujarati for to inform, then say \emph{patrakār} again
\end{itemize}

\begin{tcolorbox}[breakable,colback=teal!4,colframe=teal!35!black,title={Before you move on}]
What does \gu{પત્રકાર} mean? (\textquotedblleft{}A journalist\textquotedblright{}.) Say it once more.
\end{tcolorbox}
\section[\texorpdfstring{\gu{ઘટના} (ghaṭnā)}{ઘટના (ghaṭnā)}]{\gu{ઘટના} (ghaṭnā) — an incident; an event}

\label{lesson:GU-C312-ghatna}



\begin{quote}
Before the new one: say the Gujarati for to inform, then the Gujarati for a journalist.
\end{quote}

\subsection*{You'll want to know: \gu{ઘટના}}
\textbf{\gu{ઘટના}} — \emph{ghaṭnā} — \textquotedblleft{}an incident; an event\textquotedblright{}.

\begin{quote}
\gu{પત્રકારે} \gu{આખી} \gu{ઘટના} \gu{લખી}.
\end{quote}

\begin{quote}
\emph{patrakāre ākhī ghaṭnā lakhī.}
\end{quote}

\begin{quote}
The journalist wrote down the whole incident.
\end{quote}

A feminine noun.

\begin{cousinweb}
\textbf{\gu{ઘટના}} (\emph{ghaṭnā}) is a tatsama, Sanskrit \emph{ghaṭanā}, \textquotedblleft{}taking effect, a coming together\textquotedblright{} (Monier-Williams); in Gujarati it is something that happens, the word of reports: \textbf{\gu{આખી} \gu{ઘટના}}, the whole incident.
\end{cousinweb}

\begin{grammarlens}[title={What you've built}]
One more word for telling what happened.
\end{grammarlens}

\subsection*{Your turn}
\begin{itemize}
\raggedright
  \item \emph{Say it:} \emph{ghaṭnā}
  \item \emph{Say it:} \emph{ghaṭnā}, once more
  \item \emph{Say it:} \emph{patrakāre ākhī ghaṭnā lakhī.} — The journalist wrote down the whole incident.
  \item \emph{Recall:} say the Gujarati for to inform, then the Gujarati for a journalist, then say \emph{ghaṭnā} again
\end{itemize}

\begin{tcolorbox}[breakable,colback=teal!4,colframe=teal!35!black,title={Before you move on}]
What does \gu{ઘટના} mean? (\textquotedblleft{}An incident; an event\textquotedblright{}.) Say it once more.
\end{tcolorbox}
\section[\texorpdfstring{\gu{કિસ્સો} (kisso)}{કિસ્સો (kisso)}]{\gu{કિસ્સો} (kisso) — a case (to tell); an incident}

\label{lesson:GU-C312-kisso}



\begin{quote}
Before the new one: say the Gujarati for a journalist, then the Gujarati for an incident.
\end{quote}

\subsection*{You'll want to know: \gu{કિસ્સો}}
\textbf{\gu{કિસ્સો}} — \emph{kisso} — \textquotedblleft{}a case (to tell); an incident\textquotedblright{}.

\begin{quote}
\gu{પત્રકારે} \gu{છાપામાં} \gu{આ} \gu{કિસ્સો} \gu{લખ્યો}.
\end{quote}

\begin{quote}
\emph{patrakāre chāpāmā̃ ā kisso lakhyo.}
\end{quote}

\begin{quote}
The journalist wrote this story up in the paper.
\end{quote}

A masculine noun.

\begin{cousinweb}
\textbf{\gu{કિસ્સો}} (\emph{kisso}) is usually traced, through Persian, to Arabic \emph{qiṣṣa}, \textquotedblleft{}a story\textquotedblright{}. It is an incident told as a story: \textbf{\gu{એક} \gu{કિસ્સો}}, a case, something that happened to someone.
\end{cousinweb}

\begin{grammarlens}[title={What you've built}]
One more word for telling what happened.
\end{grammarlens}

\subsection*{Your turn}
\begin{itemize}
\raggedright
  \item \emph{Say it:} \emph{kisso}
  \item \emph{Say it:} \emph{kisso}, once more
  \item \emph{Say it:} \emph{patrakāre chāpāmā̃ ā kisso lakhyo.} — The journalist wrote this story up in the paper.
  \item \emph{Recall:} say the Gujarati for a journalist, then the Gujarati for an incident, then say \emph{kisso} again
\end{itemize}

\begin{tcolorbox}[breakable,colback=teal!4,colframe=teal!35!black,title={Before you move on}]
What does \gu{કિસ્સો} mean? (\textquotedblleft{}A case (to tell); an incident\textquotedblright{}.) Say it once more.
\end{tcolorbox}
\section[\texorpdfstring{\gu{અહેવાલ} (ahevāl)}{અહેવાલ (ahevāl)}]{\gu{અહેવાલ} (ahevāl) — a report}

\label{lesson:GU-C312-aheval}



\begin{quote}
Before the new one: say the Gujarati for an incident, then the Gujarati for a case (to tell).
\end{quote}

\subsection*{You'll want to know: \gu{અહેવાલ}}
\textbf{\gu{અહેવાલ}} — \emph{ahevāl} — \textquotedblleft{}a report\textquotedblright{}.

\begin{quote}
\gu{બીજા} \gu{દિવસે} \gu{છાપામાં} \gu{તેનો} \gu{અહેવાલ} \gu{આવ્યો}.
\end{quote}

\begin{quote}
\emph{bījā divse chāpāmā̃ teno ahevāl āvyo.}
\end{quote}

\begin{quote}
The next day his report appeared in the paper.
\end{quote}

A masculine noun.

\begin{cousinweb}
\textbf{\gu{અહેવાલ}} (\emph{ahevāl}) is usually traced, through Persian, to Arabic \emph{aḥwāl}, \textquotedblleft{}states, conditions\textquotedblright{}, the plural of the word behind \textbf{\gu{હાલત}}: an account of how things stand.
\end{cousinweb}

\begin{grammarlens}[title={What you've built}]
One more word for telling what happened.
\end{grammarlens}

\subsection*{Your turn}
\begin{itemize}
\raggedright
  \item \emph{Say it:} \emph{ahevāl}
  \item \emph{Say it:} \emph{ahevāl}, once more
  \item \emph{Say it:} \emph{bījā divse chāpāmā̃ teno ahevāl āvyo.} — The next day his report appeared in the paper.
  \item \emph{Recall:} say the Gujarati for an incident, then the Gujarati for a case (to tell), then say \emph{ahevāl} again
\end{itemize}

\begin{tcolorbox}[breakable,colback=teal!4,colframe=teal!35!black,title={Before you move on}]
What does \gu{અહેવાલ} mean? (\textquotedblleft{}A report\textquotedblright{}.) Say it once more.
\end{tcolorbox}
\section[\texorpdfstring{\gu{વર્ણવવું} (varṇavvũ)}{વર્ણવવું (varṇavvũ)}]{\gu{વર્ણવવું} (varṇavvũ) — to describe}

\label{lesson:GU-C312-varnavvu}



\begin{quote}
Before the new one: say the Gujarati for a case (to tell), then the Gujarati for a report.
\end{quote}

\subsection*{You'll want to know: \gu{વર્ણવવું}}
\textbf{\gu{વર્ણવવું}} — \emph{varṇavvũ} — \textquotedblleft{}to describe\textquotedblright{}.

\begin{quote}
\gu{મીરાએ} \gu{પત્રકારને} \gu{બધું} \gu{વર્ણવ્યું}.
\end{quote}

\begin{quote}
\emph{Mīrāe patrakārne badhũ varṇavyũ.}
\end{quote}

\begin{quote}
Mira described everything to the journalist.
\end{quote}

Past, with \gu{તેણે} (he or she …): \gu{તેણે} \gu{બધું} \gu{વર્ણવ્યું}, \emph{teṇe badhũ varṇavyũ}.

\begin{cousinweb}
\textbf{\gu{વર્ણવવું}} (\emph{varṇavvũ}) is built on Sanskrit \emph{varṇana}, \textquotedblleft{}description\textquotedblright{} (CDIAL 11340), made into a Gujarati verb: to tell what something was like.
\end{cousinweb}

\begin{grammarlens}[title={What you've built}]
One more word for telling what happened.
\end{grammarlens}

\subsection*{Your turn}
\begin{itemize}
\raggedright
  \item \emph{Say it:} \emph{varṇavvũ}
  \item \emph{Say it:} \emph{varṇavvũ}, once more
  \item \emph{Say it:} \emph{Mīrāe patrakārne badhũ varṇavyũ.} — Mira described everything to the journalist.
  \item \emph{Recall:} say the Gujarati for a case (to tell), then the Gujarati for a report, then say \emph{varṇavvũ} again
\end{itemize}

\begin{tcolorbox}[breakable,colback=teal!4,colframe=teal!35!black,title={Before you move on}]
What does \gu{વર્ણવવું} mean? (\textquotedblleft{}To describe\textquotedblright{}.) Say it once more.
\end{tcolorbox}
